\item Está trabajando como asistente de investigación para un profesor cuya área de investigación es la termodinámica. Él le señala que Daniel Fahrenheit utilizó la mejor estimación de la temperatura normal del cuerpo humano como uno de los puntos en la definición de la escala de temperatura Fahrenheit original. En la escala revisada que usamos ahora, la temperatura normal del cuerpo humano es de 98.6 °F. Su profesor propone una nueva escala en la cual la temperatura normal del cuerpo humano sería exactamente 100 °N, donde la unidad °N es un grado en la Nueva escala. La temperatura del agua helada sería 0 °N, como en la escala Celsius. Su profesor le pide que determine las siguientes temperaturas en su nueva escala: (a) cero absoluto, (b) el punto de fusión del mercurio (-37.9 °F), (c) el punto de ebullición del agua y, para la publicidad en su futura esperada conferencia de prensa, (d) la temperatura del aire más alta registrada en la superficie de la Tierra, 134.1 °F el 10 de julio de 1913, en el Valle de la Muerte, California.
